\section{Proposed Solution}
\label{sec:proposed}

GRASP organises disclosure review around the structured BRSR report. Narrative text segments and reported facts are stored with the same company and reporting-period identifiers, creating a common source for classification, review checks, and scoring. This design allows an analyst to move from a company summary to the disclosure behind a result. The approach therefore addresses the central review problem identified in Section~\ref{sec:introduction}: connecting written claims, filed evidence, and comparable company-level outputs without losing the reporting context.

The architecture separates model training from report analysis. Report-derived weak labels support training a domain-adapted BERT classifier for the nine principles. During analysis, the trained classifier organises text segments and produces principle probabilities. Word lists and sentence similarity identify potentially vague wording, while numerical and assurance rules generate separate evidence-linked review records. The Transparency Score combines four disclosure measures and reuses the vagueness detector for Clarity, which carries a weight of 20\%. Principle probabilities and numerical flags are presented alongside the score rather than included as score components. Figure~\ref{fig:architecture} shows the relationships between these functions; Section~\ref{sec:methodology} defines the calculations.

Implementation requires the collected reports, a database, pretrained language models, and computing resources for model training and report analysis. The trained classifier can then be reused by the service layer, which exposes company records, principle summaries, review flags, scores, and peer comparisons through an application programming interface (API). These outputs support examination of one company, selected companies in a sector, the complete financial-sector group, and all recorded sectors. Where two annual reports are available, the same scoring rules also support year-wise comparison. Evaluation in Section~\ref{sec:evaluation} distinguishes model agreement, language-screening performance, score comparisons, and service verification.
